\chapter{Introducción}
\label{sec:spec-intro}

Este libro es la \emph{especificación} de las superficies externas
de GraphHunter. Enumera, con contrato, todo lo que el sistema
expone al mundo: los 64 endpoints HTTP, las 55 tools MCP, la
gramática del DSL de hipótesis, la C ABI de \filepath{libgraphmatch},
los formatos de log aceptados, las variables de entorno y los
códigos de error.

\textbf{Véase también} el libro \emph{Diseño completo}
(\filepath{docs/diseno/diseno.pdf}) para entender el \emph{porqué}
detrás de cada \emph{contract} --- por ejemplo, por qué hay tantas
variables de entorno relacionadas a \emph{indexes} (heavy-logs
\#1--\#4) o cómo el separador \texttt{;} de la DSL se conecta con
el \emph{tree-walk Yannakakis enumerator}.

\section{Para qué sirve esta especificación}

Un cliente del sistema — humano o automatizado — necesita saber qué
puede pedir, qué va a recibir y qué cambios de comportamiento
pueden esperarse entre versiones. Este libro responde esas tres
preguntas en forma de contratos. Cada contrato indica:

\begin{itemize}
    \item Su \textbf{estabilidad}: \code{stable}, \code{experimental}
    o \code{deprecated}.
    \item Su \textbf{versión} (semántica) cuando aplica.
    \item Los \textbf{tipos de entrada y salida}, con un ejemplo.
    \item Los \textbf{códigos de error} que puede emitir.
\end{itemize}

\section{Qué no está aquí}

Los algoritmos, los patrones de implementación y las decisiones
arquitectónicas viven en el libro de Arquitectura. El flujo de uso
para un analista vive en el Manual. La justificación experimental
de la composición de scoring y la equivalencia DFS/SIMD vive en el
Paper.

\section{Estabilidad y versionado}

El sistema adhiere a semver con un matiz: cada superficie tiene su
propia curva. Los DTOs de la HTTP API viven en
\filepath{platform/api/src/dto/v1/} y están versionados como
\code{v1}. Las tools MCP llevan \code{version} por tool (ver ADR-004
en el libro de Arquitectura). El DSL tiene \code{v1} estable y
\code{v2-experimental} feature-gated. La C ABI lleva
\code{gm_abi_version()}.

La regla operativa: \textbf{no} se rompe silenciosamente una
superficie \code{stable}. Si una release cambia el comportamiento
visible, sube \code{version} y se documenta el cambio en el
apéndice~\ref{app:spec-changelog}.

\section{Convenciones del libro}

\begin{itemize}
    \item \textbf{Endpoints HTTP}: se listan como
    \code{METHOD /ruta} con descripción corta, payload JSON de
    request y response, códigos de error relevantes.
    \item \textbf{Tools MCP}: se listan con \code{nombre},
    \code{domain}, \code{stability}, \code{version}, input schema y
    output schema.
    \item \textbf{Funciones C}: se listan con signature, pre/post
    condiciones y ownership semántico.
    \item \textbf{Producciones BNF} para la gramática DSL.
    \item \textbf{Env vars}: nombre, tipo esperado, default,
    componente que la consume.
\end{itemize}

\section{Mapa del libro}

\begin{itemize}
    \item Cap.~\ref{sec:api-http} — HTTP API: 64 endpoints.
    \item Cap.~\ref{sec:mcp-tools} — MCP: 55 tools en 10 dominios.
    \item Cap.~\ref{sec:dsl} — DSL de hipótesis: gramática BNF y
    ejemplos.
    \item Cap.~\ref{sec:log-formats} — Formatos de log aceptados
    por los parsers.
    \item Cap.~\ref{sec:cabi} — C ABI de \filepath{libgraphmatch}.
    \item Cap.~\ref{sec:scoring} — Scoring: cinco componentes con
    dominio y fórmula.
    \item Cap.~\ref{sec:envvars} — Variables de entorno.
    \item Cap.~\ref{sec:errores} — Errores por superficie.
    \item Apéndice~\ref{app:spec-changelog} — Changelog de
    superficies.
\end{itemize}
